(1) 静态方程$-c^2\phi''+\omega_0^2\sin\phi=0$，乘$\phi'$积分得$c^2\phi'^2/2=\omega_0^2(1-\cos\phi)$，积分常数由无穷远真空边界取零。取$\ell=c/\omega_0$，分离变量得到$\phi_K(x)=4\arctan e^{(x-x_0)/\ell}$。其导数$\phi_K'=2\operatorname{sech}[(x-x_0)/\ell]/\ell$，代回方程可逐式核验。
(2) 该方程在以$c$为极限速度的时空缩放下具Lorentz型不变性。令$\gamma_v=(1-v^2/c^2)^{-1/2}$，把静态坐标替换为$\gamma_v(x-vt-x_0)$，得到$\phi_{K,v}=4\arctan\exp[\gamma_v(x-vt-x_0)/\ell]$。固定相位宽度$\ell_v=\ell/\gamma_v=\ell\sqrt{1-v^2/c^2}$，速度趋$c$时 kink 空间收缩。
(3) 静态解有$c^2\phi_x^2/2=\omega_0^2(1-\cos\phi)=2\omega_0^2\operatorname{sech}^2[(x-x_0)/\ell]$，总能量$E_0=4\omega_0^2\int\operatorname{sech}^2[(x-x_0)/\ell]dx=8c\omega_0$。移动解满足$\phi_t=-v\phi_x$且$\phi_x$增长$\gamma_v$，将其代入能量积分并换元得$E_v=\gamma_v E_0$，体现孤子整体的相对论色散。
(4) 对有限能量边界，场在$x\to\pm\infty$必须落于势的真空$\phi=2\pi n$，定义$Q=[\phi(+\infty)-\phi(-\infty)]/(2\pi)\in\mathbb Z$。时间演化若保持端点真空，$dQ/dt=[\phi_t(+\infty)-\phi_t(-\infty)]/(2\pi)=0$。本kink的$Q=1$无法连续缩成$Q=0$的小振幅线性波，除非破坏边界或引入缺陷。
